\section{AHSME 1974}
        \begin{enumerate}
        	\item (\textit{AHSME 1974}) If $x\not=0$ or $4$ and $y\not=0$ or $6$, then $\frac{2}{x}+\frac{3}{y}=\frac{1}{2}$ is equivalent to


 $\mathrm{(A)\ } 4x+3y=xy \qquad \mathrm{(B) \ }y=\frac{4x}{6-y} \qquad \mathrm{(C) \  } \frac{x}{2}+\frac{y}{3}=2 \qquad$
$\mathrm{(D) \  } \frac{4y}{y-6}=x \qquad \mathrm{(E) \  }\text{none of these}$


	\item (\textit{AHSME 1974}) Let $x_1$ and $x_2$ be such that $x_1\not=x_2$ and $3x_i^2-hx_i=b$, $i=1, 2$. Then $x_1+x_2$ equals


 $\mathrm{(A)\ } -\frac{h}{3} \qquad \mathrm{(B) \ }\frac{h}{3} \qquad \mathrm{(C) \  } \frac{b}{3} \qquad \mathrm{(D) \  } 2b \qquad \mathrm{(E) \  }-\frac{b}{3}$


	\item (\textit{AHSME 1974}) The coefficient of $x^7$ in the polynomial expansion of 


 
 \[(1+2x-x^2)^4\]


 is


 $\mathrm{(A)\ } -8 \qquad \mathrm{(B) \ }12 \qquad \mathrm{(C) \  } 6 \qquad \mathrm{(D) \  } -12 \qquad \mathrm{(E) \  }\text{none of these}$


	\item (\textit{AHSME 1974}) What is the remainder when $x^{51}+51$ is divided by $x+1$?


 $\mathrm{(A)\ } 0 \qquad \mathrm{(B) \ }1 \qquad \mathrm{(C) \  } 49 \qquad \mathrm{(D) \  } 50 \qquad \mathrm{(E) \  }51$


	\item (\textit{AHSME 1974}) Given a quadrilateral $ABCD$ inscribed in a circle with side $AB$ extended beyond $B$ to point $E$, if $\measuredangle BAD=92^\circ$ and $\measuredangle ADC=68^\circ$, find $\measuredangle EBC$.


 $\mathrm{(A)\ } 66^\circ \qquad \mathrm{(B) \ }68^\circ \qquad \mathrm{(C) \  } 70^\circ \qquad \mathrm{(D) \  } 88^\circ \qquad \mathrm{(E) \  }92^\circ$


	\item (\textit{AHSME 1974}) For positive real numbers $x$ and $y$ define $x*y=\frac{x\cdot y}{x+y}$' then


 $\mathrm{(A)\ } \text{``*" is commutative but not associative} \qquad$
$\mathrm{(B) \ }\text{``*" is associative but not commutative} \qquad$
$\mathrm{(C) \  } \text{``*" is neither commutative nor associative} \qquad$
$\mathrm{(D) \  } \text{``*" is commutative and associative} \qquad$
$\mathrm{(E) \  }\text{none of these} \qquad$


	\item (\textit{AHSME 1974}) A town's population increased by $1,200$ people, and then this new population decreased by $11\%$. The town now had $32$ less people than it did before the $1,200$ increase. What is the original population?


 $\mathrm{(A)\ } 1,200 \qquad \mathrm{(B) \ }11,200 \qquad \mathrm{(C) \  } 9,968 \qquad \mathrm{(D) \  } 10,000 \qquad \mathrm{(E) \  }\text{none of these}$


	\item (\textit{AHSME 1974}) What is the smallest prime number dividing the sum $3^{11}+5^{13}$?


 $\mathrm{(A)\ } 2 \qquad \mathrm{(B) \ }3 \qquad \mathrm{(C) \  } 5 \qquad \mathrm{(D) \  } 3^{11}+5^{13} \qquad \mathrm{(E) \  }\text{none of these}$


	\item (\textit{AHSME 1974}) The integers greater than one are arranged in five columns as follows:


 
 \[\begin{tabular}{c c c c c}\ & 2 & 3 & 4 & 5\\ 9 & 8 & 7 & 6 &\ \\ \ & 10 & 11 & 12 & 13\\ 17 & 16 & 15 & 14 &\ \\ \ & . & . & . & .\\ \end{tabular}\]


 (Four consecutive integers appear in each row; in the first, third and other odd numbered rows, the integers appear in the last four columns and increase from left to right; in the second, fourth and other even numbered rows, the integers appear in the first four columns and increase from right to left.)


 In which column will the number $1,000$ fall?


 $\mathrm{(A)\ } \text{first} \qquad \mathrm{(B) \ }\text{second} \qquad \mathrm{(C) \  } \text{third} \qquad \mathrm{(D) \  } \text{fourth} \qquad \mathrm{(E) \  }\text{fifth}$


	\item (\textit{AHSME 1974}) What is the smallest integral value of $k$ such that 


 
 \[2x(kx-4)-x^2+6=0\]


 has no real roots?


 $\mathrm{(A)\ } -1 \qquad \mathrm{(B) \ }2 \qquad \mathrm{(C) \  } 3 \qquad \mathrm{(D) \  } 4 \qquad \mathrm{(E) \  }5$


	\item (\textit{AHSME 1974}) If $(a, b)$ and $(c, d)$ are two points on the line whose equation is $y=mx+k$, then the distance between $(a, b)$ and $(c, d)$, in terms of $a, c,$ and $m$ is 


 $\mathrm{(A)\ } |a-c|\sqrt{1+m^2} \qquad \mathrm{(B) \ }|a+c|\sqrt{1+m^2} \qquad \mathrm{(C) \  } \frac{|a-c|}{\sqrt{1+m^2}} \qquad$
$\mathrm{(D) \  } |a-c|(1+m^2) \qquad \mathrm{(E) \  }|a-c|\,|m|$


	\item (\textit{AHSME 1974}) If $g(x)=1-x^2$ and $f(g(x))=\frac{1-x^2}{x^2}$ when $x\not=0$, then $f(1/2)$ equals 


 $\mathrm{(A)\ } 3/4 \qquad \mathrm{(B) \ }1 \qquad \mathrm{(C) \  } 3 \qquad \mathrm{(D) \  } \sqrt{2}/2 \qquad \mathrm{(E) \  }\sqrt{2}$


	\item (\textit{AHSME 1974}) Which of the following is equivalent to "If P is true, then Q is false."?


 $\mathrm{(A)\ } \text{``P is true or Q is false."} \qquad$
$\mathrm{(B) \ }\text{``If Q is false then P is true."} \qquad$
$\mathrm{(C) \  } \text{``If P is false then Q is true."} \qquad$
$\mathrm{(D) \  } \text{``If Q is true then P is false."} \qquad$
$\mathrm{(E) \  }\text{``If Q is true then P is true."} \qquad$


	\item (\textit{AHSME 1974}) Which statement is correct?


 $\mathrm{(A)\ } \text{If } x<0, \text{then } x^2>x. \qquad \mathrm{(B) \ } \text{If } x^2>0, \text{then } x>0.$
$\qquad \mathrm{(C) \  } \text{If } x^2>x, \text{then } x>0. \qquad \mathrm{(D) \  } \text{If } x^2>x, \text{then } x<0.$
$\qquad \mathrm{(E) \  }\text{If } x<1, \text{then } x^2<x.$


	\item (\textit{AHSME 1974}) If $x<-2$, then $|1-|1+x||$ equals 


 $\mathrm{(A)\ } 2+x \qquad \mathrm{(B) \ }-2-x \qquad \mathrm{(C) \  } x \qquad \mathrm{(D) \  } -x \qquad \mathrm{(E) \  }-2$


	\item (\textit{AHSME 1974}) A circle of radius $r$ is inscribed in a right isosceles triangle, and a circle of radius $R$ is circumscribed about the triangle. Then $R/r$ equals


 $\mathrm{(A)\ } 1+\sqrt{2} \qquad \mathrm{(B) \ }\frac{2+\sqrt{2}}{2} \qquad \mathrm{(C) \  } \frac{\sqrt{2}-1}{2} \qquad \mathrm{(D) \  } \frac{1+\sqrt{2}}{2} \qquad \mathrm{(E) \  }2(2-\sqrt{2})$


	\item (\textit{AHSME 1974}) If $i^2=-1$, then $(1+i)^{20}-(1-i)^{20}$ equals


 $\mathrm{(A)\ } -1024 \qquad \mathrm{(B) \ }-1024i \qquad \mathrm{(C) \  } 0 \qquad \mathrm{(D) \  } 1024 \qquad \mathrm{(E) \  }1024i$


	\item (\textit{AHSME 1974}) If $\log_8{3}=p$ and $\log_3{5}=q$, then, in terms of $p$ and $q$, $\log_{10}{5}$ equals


 $\mathrm{(A)\ } pq \qquad \mathrm{(B) \ }\frac{3p+q}{5} \qquad \mathrm{(C) \  } \frac{1+3pq}{p+q} \qquad \mathrm{(D) \  } \frac{3pq}{1+3pq} \qquad \mathrm{(E) \  }p^2+q^2$


	\item (\textit{AHSME 1974}) In the adjoining figure $ABCD$ is a square and $CMN$ is an equilateral triangle. If the area of $ABCD$ is one square inch, then the area of $CMN$ in square inches is 


 
                    \begin{figure}[H]
                        \centering
                        \includegraphics[width=0.5\textwidth]{Problem/AHSME_I_1974_19_Problem_Figure_1.png}
                    \end{figure}
                    
$\mathrm{(A)\ } 2\sqrt{3}-3 \qquad \mathrm{(B) \ }1-\frac{\sqrt{3}}{3} \qquad \mathrm{(C) \  } \frac{\sqrt{3}}{4} \qquad \mathrm{(D) \  } \frac{\sqrt{2}}{3} \qquad \mathrm{(E) \  }4-2\sqrt{3}$


	\item (\textit{AHSME 1974}) Let


 $T=\frac{1}{3-\sqrt{8}}-\frac{1}{\sqrt{8}-\sqrt{7}}+\frac{1}{\sqrt{7}-\sqrt{6}}-\frac{1}{\sqrt{6}-\sqrt{5}}+\frac{1}{\sqrt{5}-2}.$


 Then


 $\mathrm{(A)\ } T<1 \qquad \mathrm{(B) \ }T=1 \qquad \mathrm{(C) \  } 1<T<2 \qquad \mathrm{(D) \  } T>2 \qquad$
$\mathrm{(E) \  }T=\frac{1}{(3-\sqrt{8})(\sqrt{8}-\sqrt{7})(\sqrt{7}-\sqrt{6})(\sqrt{6}-\sqrt{5})(\sqrt{5}-2)}$


	\item (\textit{AHSME 1974}) In a geometric series of positive terms the difference between the fifth and fourth terms is $576$, and the difference between the second and first terms is $9$. What is the sum of the first five terms of this series?


 $\mathrm{(A)\ } 1061 \qquad \mathrm{(B) \ }1023 \qquad \mathrm{(C) \  } 1024 \qquad \mathrm{(D) \  } 768 \qquad \mathrm{(E) \  }\text{none of these}$


	\item (\textit{AHSME 1974}) The minimum of $\sin\frac{A}{2}-\sqrt{3}\cos\frac{A}{2}$ is attained when $A$ is 


 $\mathrm{(A)\ } -180^\circ \qquad \mathrm{(B) \ }60^\circ \qquad \mathrm{(C) \  } 120^\circ \qquad \mathrm{(D) \  } 0^\circ \qquad \mathrm{(E) \  }\text{none of these}$


	\item (\textit{AHSME 1974}) In the adjoining figure $TP$ and $T'Q$ are parallel tangents to a circle of radius $r$, with $T$ and $T'$ the points of tangency. $PT''Q$ is a third tangent with $T'''$ as a point of tangency. If $TP=4$ and $T'Q=9$ then $r$ is


 
                    \begin{figure}[H]
                        \centering
                        \includegraphics[width=0.5\textwidth]{Problem/AHSME_I_1974_23_Problem_Figure_1.png}
                    \end{figure}
                    
$\mathrm{(A)\ } 25/6 \qquad \mathrm{(B) \ } 6 \qquad \mathrm{(C) \  } 25/4 \qquad$
$\mathrm{(D) \  } \text{a number other than }25/6, 6, 25/4 \qquad$
$\mathrm{(E) \  }\text{not determinable from the given information}$


	\item (\textit{AHSME 1974}) A fair die is rolled six times. The probability of rolling at least a five at least five times is 


 $\mathrm{(A)\ } \frac{13}{729} \qquad \mathrm{(B) \ }\frac{12}{729} \qquad \mathrm{(C) \  } \frac{2}{729} \qquad \mathrm{(D) \  } \frac{3}{729} \qquad \mathrm{(E) \  }\text{none of these}$


	\item (\textit{AHSME 1974}) In parallelogram $ABCD$ of the accompanying diagram, line $DP$ is drawn bisecting $BC$ at $N$ and meeting $AB$ (extended) at $P$. From vertex $C$, line $CQ$ is drawn bisecting side $AD$ at $M$ and meeting $AB$ (extended) at $Q$. Lines $DP$ and $CQ$ meet at $O$. If the area of parallelogram $ABCD$ is $k$, then the area of the triangle $QPO$ is equal to 


 
                    \begin{figure}[H]
                        \centering
                        \includegraphics[width=0.5\textwidth]{Problem/AHSME_I_1974_25_Problem_Figure_1.png}
                    \end{figure}
                    
$\mathrm{(A)\ } k \qquad \mathrm{(B) \ }\frac{6k}{5} \qquad \mathrm{(C) \  } \frac{9k}{8} \qquad \mathrm{(D) \  } \frac{5k}{4} \qquad \mathrm{(E) \  }2k$


	\item (\textit{AHSME 1974}) The number of distinct positive integral divisors of $(30)^4$ excluding $1$ and $(30)^4$ is 


 $\mathrm{(A)\ } 100 \qquad \mathrm{(B) \ }125 \qquad \mathrm{(C) \  } 123 \qquad \mathrm{(D) \  } 30 \qquad \mathrm{(E) \  }\text{none of these}$


	\item (\textit{AHSME 1974}) If $f(x)=3x+2$ for all real $x$, then the statement:
"$|f(x)+4|<a$ whenever $|x+2|<b$ and $a>0$ and $b>0$"
is true when


 $\mathrm{(A)}\ b\le a/3\qquad\mathrm{(B)}\ b > a/3\qquad\mathrm{(C)}\ a\le b/3\qquad\mathrm{(D)}\ a > b/3\\ \qquad\mathrm{(E)}\ \text{The statement is never true.}$


	\item (\textit{AHSME 1974}) Which of the following is satisfied by all numbers $x$ of the form


 
 \[x=\frac{a_1}{3}+\frac{a_2}{3^2}+\cdots+\frac{a_{25}}{3^{25}}\]


 where $a_1$ is $0$ or $2$, $a_2$ is $0$ or $2$,...,$a_{25}$ is $0$ or $2$?


 $\mathrm{(A)\ } 0\le x<1/3 \qquad \mathrm{(B) \ } 1/3\le x<2/3 \qquad \mathrm{(C) \  } 2/3\le x<1 \qquad$
$\mathrm{(D) \  } 0\le x<1/3 \text{ or }2/3\le x<1 \qquad \mathrm{(E) \  }1/2\le x\le 3/4$


	\item (\textit{AHSME 1974}) For $p=1, 2, \cdots, 10$ let $S_p$ be the sum of the first $40$ terms of the arithmetic progression whose first term is $p$ and whose common difference is $2p-1$; then $S_1+S_2+\cdots+S_{10}$ is


 $\mathrm{(A)\ } 80000 \qquad \mathrm{(B) \ }80200 \qquad \mathrm{(C) \  } 80400 \qquad \mathrm{(D) \  } 80600 \qquad \mathrm{(E) \  }80800$


	\item (\textit{AHSME 1974}) A line segment is divided so that the lesser part is to the greater part as the greater part is to the whole. If $R$ is the ratio of the lesser part to the greater part, then the value of 


 
 \[R^{[R^{(R^2+R^{-1})}+R^{-1}]}+R^{-1}\]


 is


 $\mathrm{(A)\ } 2 \qquad \mathrm{(B) \ }2R \qquad \mathrm{(C) \  } R^{-1} \qquad \mathrm{(D) \  } 2+R^{-1} \qquad \mathrm{(E) \  }2+R$


\end{enumerate}